%% title: Accuracy saturates; the three axes disagree
%% target: Sections/Results.tex
%% anchor: Shift lowers the logits of known species towards those of novel ones, so a species close to a known one passes as known: novelty detection is not robust to the domain, and a species-level answer on a trap image means ``this species, or a close relative the model never saw''.
%% action: after
%% why: The motivating table: models within 1.4 points on the test fold span 43 points on trap images, and the axes rank them differently, so the test fold cannot select a model for deployment.
%% source: paper/DRAFT.md §4.0, §4.2, §4.10; data/rescore/probe-grouped.csv and probe-ho-grouped.csv (base-probef 0.2875 / 0.2394; B1 0.4242 / 0.3561; B8 0.6472 / 0.4520; P5 0.6223 / 0.4974; T1rep 0.7232 / 0.3122); journal/2026-10-22-rescoring-in-the-mini-metrics-convention.md (in-distribution)
%% check: Paste after step 14. | Recipe-L models, before the rerun; the rerun's head grid replaces the rows. | The model names B1, B8, P5, T1rep need defining once in the paper (here, or in a model table in the methods).
%% status: ready
\subsection{Accuracy on the test fold does not predict accuracy under shift}\label{sec:results_axes}
\Cref{tab:axes} scores five recipe-L models on every axis.
On the test fold they lie within 1.4 points; on the trap probe set they span 43 points, and on its held-out species 26.
The most accurate model we trained on the test fold (ConvNeXtV2-L at 320 pixels, 0.932) is 1.8 points above the baseline there, and self-training on trap images, worth 21 points on the probe set (\cref{sec:results_selftraining}), changes the test fold by 0.04.
The axes also rank models differently: real trap labels give the best model on the probe set and the worst adapted model on the species those labels never covered.

\begin{table}[htb]
    \centering
    \caption{Species macro-$\text{F}_1$ on the three axes (recipe L). Test fold: zero threshold. Probe and held-out species: calibrated on 10\,\% of the (trap, night) groups. Open set: AUROC with the model's best rule, rare-species benchmark.}
    \label{tab:axes}
    \begin{tabular}{lcccc}
        \toprule
        model & test fold & probe & held-out species & open set \\
        \midrule
        F, EfficientNetV2-S, no trap images & 0.9135 & 0.288 & 0.239 & 0.899 \\
        B1: + margin and trap-like augmentation & 0.8999 & 0.424 & 0.356 & -- \\
        B8: ConvNeXt-L, self-training, no oversampling & 0.9060 & 0.647 & 0.452 & 0.915 \\
        P5: BioCLIP-2 fine-tuned, adapted & 0.9113 & 0.622 & 0.497 & 0.916 \\
        T1rep: B1 + real trap labels, 11 copies & -- & 0.723 & 0.312 & -- \\
        \bottomrule
    \end{tabular}
\end{table}

Three axes are therefore needed, though the two strongest general models, B8 and P5, are ahead of the baseline on all three and close to each other.
What none of them shows is how a model abstains (\cref{sec:results_backoff}).
